%% kapa: exam=1 round=32 year=2021 period=1 form=A subject=real_estate subject_ko=부동산학원론 seq=17 number=97
%% source: Q-Net 감정평가사 기출문제 (cst00309) artlSeq=5208397 file=extracted/1008/5208397/A형/2021년도 제32회 감정평가사 1차 1교시 A형.pdf
%% conversion: kordoc 4.24.51 -> tools/gen1.py (auto); review=proofread-sonnet; flags=none
\begin{question}{97}
포트폴리오 이론에 관한 설명으로 옳지 않은 것은?
\begin{choices}
  \item 부동산투자에 수반되는 총위험은 체계적 위험과 비체계적 위험을 합한 것으로, 포트폴리오를 구성함으로써 제거될 수 있는 위험은 비체계적 위험이다.
  \item 포트폴리오를 구성하는 자산들의 수익률 간 상관계수가 1인 경우에는 포트폴리오를 구성한다고 하더라도 위험은 감소되지 않는다.
  \item 효율적 프론티어(efficient frontier)는 모든 위험수준에서 최대의 기대수익률을 올릴 수 있는 포트폴리오의 집합을 연결한 선이다.
  \item 무위험자산이 없는 경우의 최적 포트폴리오는 효율적 프론티어(efficient frontier)와 투자자의 무차별곡선이 접하는 점에서 결정되는데, 투자자가 위험선호형일 경우 최적 포트폴리오는 위험기피형에 비해 저위험-고수익 포트폴리오가 된다.
  \item 위험자산으로만 구성된 포트폴리오와 무위험자산을 결합할 때 얻게 되는 직선의 기울기가 커질수록 기대초과수익률(위험프리미엄)이 커진다.
\end{choices}
\end{question}
\answer{4}
